% Appendix E -- trained backward policies: flow regularization in spectral variables.
\section{\appendixEtitle}\label{app:pbtrain}

Everything in Section~\ref{sec:cv_divergence} freezes the backward policy; the practice
recommended by \citet{morozov2025revisiting} trains it. Their objective is the log-scale
detailed-balance loss (their Eq.~4) with a learnable $\mathcal P_B$, plus state-flow
regularization (their Eq.~12),
\begin{equation}\label{equ:hse_objective}
    \mathcal L_{\lamreg}(\mathcal F, \mathcal P_B)
    \;=\; \sum_{s\to s'} \widehat\nu(s\to s')
    \Bigl(\log\tfrac{\mathcal F(s)\mathcal P_F(s'|s)}{\mathcal F(s')\mathcal P_B(s|s')}\Bigr)^2
    \;+\; \lamreg\sum_{s}\mathcal F(s)
\end{equation}
(their regularizer is written per transition; summed over a batch it is the total state
flow, which is the form every statement below uses), with reward matching enforced by
the substitution
$\mathcal F(s_f)\mathcal P_B(x|s_f)=\mathcal R(x)$. That substitution does two things at
once in the loop-closure frame of this paper: it pins the total scale ($\mathcal F(s_f) =
\sum_x\mathcal R(x) = Z$ on the balanced set, so the scale ray of the frozen-backward
family --- Theorem~\ref{theo:universality_graphs}(2), the ray whose uniqueness-up-to-scale
Proposition~\ref{prop:frozen_unstable} exploits --- is removed by a boundary condition,
not by a constraint one must add), and it pins the \emph{sink row} of the backward chain to the
normalized target --- exactly the convention
$\widehat{\mathcal P}_B(s_f\to\cdot)=\target/t$ under which
Proposition~\ref{prop:morozov_rate} defines $N$, $\bar\sigma$, $\sigma_*$. Their
Proposition~3.7 then parameterizes the balanced set by the \emph{interior rows of
$\mathcal P_B$ alone}. This appendix translates \eqref{equ:hse_objective} into the
spectral variables of Section~\ref{sec:cv_divergence}: what the $\mathcal P_B$-training
does (a slow drift along the balanced manifold, descending expected trajectory length),
where it goes (the optimal-transport endpoint of \citet{maksimov2026ot}), and what it
costs (every constant of the training-speed theory, at a computable exchange rate).

\subsection{The joint landscape, and the dictionary}\label{sec:joint}

\begin{lemma}[balanced points are critical in all variables]\label{lem:balanced_min}
Let $\mathcal L\geq0$ be $C^1$ in $(\mathcal F,\mathcal P_B)$ and vanish exactly on the
balanced set. Then every balanced point is a global minimum, so
$\nabla_{(\mathcal F,\mathcal P_B)}\mathcal L=0$ there --- including the
$\mathcal P_B$-gradient. If moreover $\mathcal L=\sum\widehat\nu\, g(r)$ with
$g(1)=g'(1)=0$, $g''(1)>0$, the joint Hessian at a balanced point is Gauss--Newton,
$H_{\mathrm{full}} = g''(1)\, J^\dagger M_{\widehat\nu} J$ with $J$ the joint Jacobian of
the residual $r$, and $\ker H_{\mathrm{full}}$ is the tangent space of the balanced
manifold.
\end{lemma}

\begin{proof}
The first claim is the definition of a global minimum of a nonnegative function. For the
second, $\nabla^2[g(r)] = g''(r)\nabla r\otimes\nabla r + g'(r)\nabla^2 r$; at $r\equiv1$
the second term dies with $g'(1)=0$.
\end{proof}

So at balance the only surviving force in \eqref{equ:hse_objective} is the regularizer:
$\mathcal P_B$-training near the balanced manifold is not ``the loss adjusting the
backward policy'' --- the loss is indifferent there --- it is the regularizer
transporting the flow \emph{along} the manifold. What the regularizer descends is
identified exactly by the following dictionary entry.

\begin{lemma}[their Eq.~10 is the occupation identity, summed]\label{lem:eq10_occupation}
On the loop closure with sink row $\target/t$, let $\mathcal F$ be the balanced flow of an
interior-positive $\mathcal P_B$ with $\mathcal F(s_f)=Z$. Then $\mathcal F(s) = Z\,N(s)$
for every state, hence
\begin{equation}\label{equ:eq10}
    \sum_{s\notin\{s_0,s_f\}}\frac{\mathcal F(s)}{\mathcal F(s_f)} \;=\; \bar\sigma,
    \qquad
    \sum_{s}\frac{\mathcal F(s)}{\mathcal F(s_f)} \;=\; 2+\bar\sigma,
    \qquad
    \mathcal F(\Statespace) \;=\; (2+\bar\sigma)\,Z .
\end{equation}
The first display is Eq.~(10) of \citet{morozov2025revisiting} (expected trajectory
length $=$ normalized total interior flow) re-derived in the variables of
Proposition~\ref{prop:morozov_rate}; in particular the state-flow regularizer of
\eqref{equ:hse_objective} equals, on the balanced manifold,
$\lamreg{}Z\,(2+\bar\sigma(\mathcal P_B))$ --- up to an additive constant,
$\lamreg{}Z\,\bar\sigma(\mathcal P_B)$, so \textbf{flow regularization is
gradient descent on the expected backward-trajectory length}.
\end{lemma}

\begin{proof}
The balanced flow of a frozen backward policy is the invariant-measure ray
(Theorem~\ref{theo:universality_graphs}, part (2)): $\mathcal F = c\,\lambda$. The
occupation identity $\lambda = N/(2+\bar\sigma)$
(Proposition~\ref{prop:morozov_rate}(1)) with $\lambda(s_f)=1/(2+\bar\sigma)$ gives
$c = Z(2+\bar\sigma)$ from $\mathcal F(s_f)=Z$, hence $\mathcal F = Z N$. Sum, using
$N(s_0)=N(s_f)=1$.
\end{proof}

The paper's training-speed theorem is stated as a two-phase convergence result, not as
a gradient-domination inequality; the drift analysis below needs the latter, so we
record it as a lemma. It is the coercivity of Proposition~\ref{prop:morozov_rate} read
on the Hessian.

\begin{lemma}[local Polyak--{\L}ojasiewicz inequality]\label{lem:local_pl}
Fix an interior backward policy $p$ with occupation measure $\lambda(p)$ and training
density $w\geq w_{\min}>0$, and consider the frozen-$p$ loss at its quadratic
approximation around the balanced ray, $\mathcal L=\tfrac12\langle h,Hh\rangle$ with
$H=g''(1)A^\dagger M_wA$ on $L^2(\lambda(p))$ (Theorem~\ref{theo:gd_diffusion}). Then,
on the orthogonal complement of $\ker H$,
\begin{equation*}
    \|\nabla\mathcal L\|^2_{L^2(\lambda(p))}\;\geq\;2\,\varrho_\sigma(p)\,\mathcal L .
\end{equation*}
The constant lives in the $L^2(\lambda(p))$ metric, which moves with $p$: along the
drift studied below it degenerates exactly as $\lambda_{\min}(p)\to0$
(Proposition~\ref{prop:constants_path}).
\end{lemma}

\begin{proof}
$\nabla\mathcal L=Hh$ and, $H$ being positive semi-definite,
$\|Hh\|^2\geq\zeta^+_{\min}(H)\,\langle h,Hh\rangle=2\zeta^+_{\min}(H)\,\mathcal L$
on $(\ker H)^\perp$. The coercivity of Proposition~\ref{prop:morozov_rate}(2) bounds the
gap: $\langle h,Hh\rangle=g''(1)\langle Ah,wAh\rangle\geq
g''(1)w_{\min}\|Ah\|^2\geq g''(1)w_{\min}\|h-\Pi h\|^2/\widehat B_\sigma^{\,2}$, so
$\zeta^+_{\min}(H)\geq g''(1)w_{\min}/\widehat B_\sigma^{\,2}=\varrho_\sigma(p)$.
\end{proof}

\begin{proposition}[quasi-static tracking, and the shadow drift]\label{prop:two_timescale}
Consider the gradient flow of $\mathcal L + \lamreg R$, $R := \sum_s
\mathcal F(s)$, on a region where (i) the frozen-$\mathcal P_B$ loss satisfies the local
Polyak--{\L}ojasiewicz inequality of Lemma~\ref{lem:local_pl} with constant
$2\varrho(p)\geq2\varrho_{\min}>0$ at every backward policy $p$ visited, and (ii)
$C_R:=\sup\|\nabla R\|<\infty$. Then
\begin{equation}\label{equ:tracking}
    \mathcal L_t \;\leq\; e^{-\varrho_{\min} t}\,\mathcal L_0
    \;+\; \frac{\lamreg^2\,C_R^2}{2\,\varrho_{\min}} ,
\end{equation}
i.e.\ the trajectory tracks the balanced manifold at distance $O(\lamreg^2)$
in loss. Its slow component follows, to $O(\lamreg^2)$, the shadow ODE
$\dot p = -\lamreg\,P_{\mathrm{tan}}\nabla R = -\lamreg Z\,
P_{\mathrm{tan}}\nabla\bar\sigma(p)$: projected gradient descent on the expected
trajectory length.
\end{proposition}

\begin{proof}
Enlarging the variable space only adds nonnegative terms to the gradient norm, so
pointwise in $p_t$, $\|\nabla_{(\mathcal F,p)}\mathcal L\|^2 \geq
\|\nabla_{\mathcal F}\mathcal L\|^2 \geq 2\varrho(p_t)\mathcal L$ --- the frozen
inequality applied at the instantaneous kernel; no transverse-Hessian decomposition is
needed and the cross-terms of the joint Hessian never enter. Then
$\dot{\mathcal L} = -\langle\nabla\mathcal L, \nabla\mathcal L +
\lamreg\nabla R\rangle \leq -\tfrac12\|\nabla\mathcal L\|^2 +
\tfrac12\lamreg^2\|\nabla R\|^2 \leq -\varrho_{\min}\mathcal L +
\tfrac12\lamreg^2C_R^2$ by Young; Gr\"onwall gives \eqref{equ:tracking}.
The shadow statement is standard singular perturbation for gradient flows on a compact
interior region (the fast variable is slaved at rate $\varrho_{\min}$, the slow one moves
at $O(\lamreg)$); it is used at the level of gradient flow only ---
discrete GD adds the usual $O(\varepsilon)$ terms and is not re-derived here. The
identification $P_{\mathrm{tan}}\nabla R = Z P_{\mathrm{tan}}\nabla\bar\sigma$ is
Lemma~\ref{lem:eq10_occupation} restricted to the manifold.
\end{proof}

\subsection{Where the drift goes: the OT endpoint}\label{sec:drift_target}

\begin{lemma}[normalization reconciliation]\label{lem:normalization}
With reward-pinned terminal flows and total scale $Z$ fixed, minimizing total interior
flow over balanced flows at a \emph{fixed} initial row $\mathcal F(s_0\to u)=L(u)$ is the
linear program of \citet[Eq.~10]{maksimov2026ot}, hence by their Theorem~3.2 a Kantorovich
optimal-transport problem between $L$ and $\mathcal R$ with graph shortest-path costs.
When the initial row is learned (the setting of \eqref{equ:hse_objective}), the relevant
relaxation is $\inf_{L:\,\sum L=Z}\mathrm{OT}(L,\mathcal R)$ --- transport with a free
source marginal, still an LP, in which each reward unit is routed from its nearest
admissible entry.
\end{lemma}

\begin{proof}
A balanced flow with reward-pinned terminal edges and initial row $L$ satisfies exactly
the constraint set of \citet[Eq.~10]{maksimov2026ot}: conservation at interior states,
$\mathcal F(x\to s_f)=\mathcal R(x)$, $\mathcal F(s_0\to u)=L(u)$, non-negative edge
flows; conversely, a feasible point with everywhere-positive throughput is such a flow,
and a general feasible point lies in the \emph{closure} of the family (zero-throughput
states have no defined backward row) --- so the LP's minimum is the infimum over the
closure, the object Proposition~\ref{prop:drift_target}(ii) speaks about. The objective
--- total interior flow --- is the same on both sides, so the optimization problems
coincide and their Theorem~3.2 applies (stated through their reduced problem, Eq.~11,
under their Assumption~3.1). Freeing $L$ at fixed total $\sum L=Z$ replaces the value
by its infimum over source marginals, an LP in $(L,\Pi)$ jointly.
\end{proof}

\begin{proposition}[the drift target]\label{prop:drift_target}
(i) Along the shadow ODE of Proposition~\ref{prop:two_timescale}, $\bar\sigma(p_t)$ is
non-increasing. (ii) The global minimizers of total flow over the closure of the balanced
family are the minimum-flow solutions, supported on the tight (shortest-path) subgraph of
the optimal potentials \citep[Thms~3.2--3.3]{maksimov2026ot}. (iii) Whether the drift
\emph{converges} to a global minimizer is open --- $\bar\sigma(p)$ is a
stochastic-shortest-path policy-optimization landscape, and no claim is made about its
critical points beyond (i).
\end{proposition}

\begin{proof}
(i) is the definition of (projected) gradient descent on $Z\bar\sigma$. (ii) is
Lemma~\ref{lem:normalization} together with the cited theorems; \S\ref{sec:num_pb}
verifies the identification on enumerated instances with an independent transport LP and
a zero-reduced-cost dual certificate.
\end{proof}

\subsection{The spectral price}\label{sec:price}

The endpoint of the drift is a boundary point of the policy polytope: the min-flow
support is generically a strict subgraph, and a backward policy supported on it is
(partly) deterministic. Every constant of Section~\ref{sec:cv_divergence} is a functional
of $\mathcal P_B$; here is what each does along the way. Write $p_t$ for an interior path
with $p_t\to p_1$, $t\to1$.

\begin{proposition}[the constants along the path to minimum flow]\label{prop:constants_path}
Suppose $p_1$'s flow has a support omitting at least one \emph{state}, while every state
still reaches $s_0$ under $p_1$ (no trap). Then, as $t\to1$:
\begin{enumerate}
    \item $N_{\min}(p_t)\to0$ and $\lambda_{\min}(p_t)\to0$: the occupation of every
    off-support state vanishes;
    \item consequently $\widehat B_\sigma = \sigma_*/\lambda_{\min}^{1/2}\to\infty$ and
    $\varrho_\sigma = g''(1)w_{\min}\lambda_{\min}/\sigma_*^2\to0$: \emph{the certified
    rate dies with the off-support occupation};
    \item for any fixed full-support $\nutrain$, $\|w\|_{L^\infty} =
    \|d\nutrain/d\lambda\|_{L^\infty}\to\infty$, so the step cap
    $\varepsilon\leq(4g''(1)\|w\|_{L^\infty})^{-1}$ of
    Theorem~\ref{theo:db_stable_frozen} collapses as well;
    \item $\sigma_*$ and $\bar\sigma$ stay \emph{bounded}, converging to the hitting
    times of the limit policy. The divergence is an occupation phenomenon, not a
    hitting-time one.
\end{enumerate}
\end{proposition}

\begin{proof}
(1) $N$ is the expected-visit vector of the chain killed at $s_0$, started from the sink
row: the solve $(I-Q_p)^\top N = \widehat{\mathcal P}_B(s_f\to\cdot)$ on
$\Statespace\setminus\{s_0\}$. Under the no-trap hypothesis $I-Q_{p_1}$ is invertible
(every state reaches $s_0$), so $N$ is continuous in $p$ on a neighborhood of $p_1$.
Under $p_1$ an excursion starts on the support (the sink row is reward-pinned) and moves
only along the support, so $N(p_1)(x)=0$ for every off-support state; by hypothesis at
least one such state exists, and continuity gives $N_{\min}(p_t)\to0$. Then
$\lambda_{\min}=N_{\min}/(2+\bar\sigma)\to0$ using (4). (2) is substitution into the
definitions of $\widehat B_\sigma$ and $\varrho_\sigma$. (3) At an off-support state
$\nutrain\geq\nu_{\min}>0$ while $\lambda\to0$, so
$\|w\|_{L^\infty}\geq\nu_{\min}/\lambda_{\min}\to\infty$. (4) The hitting times solve
$(I-Q^{s_0}_p)h=\mathbf 1$, again invertible at $p_1$ by the no-trap hypothesis; by
continuity $\sigma_*(p_t)$ and $\bar\sigma(p_t)$ converge to the finite hitting times of
$p_1$.
\end{proof}

\begin{remark}[replay repairs the cap, not the rate]\label{rem:replay_rescue}
The adaptive buffer $\nutrain = \lambda(p_t)$ of Appendix~\ref{app:replay} keeps
$w\equiv1$, so the step cap survives the whole path. The rate does not:
$\varrho_\sigma = g''(1)\lambda_{\min}/\sigma_*^2\to0$ regardless. The speed loss under
flow regularization is intrinsic --- it \emph{is} the occupation mass leaving the
off-geodesic states --- and no training measure can buy it back within the guarantees of
Section~\ref{sec:cv_divergence}. (Do not read this as ``the exact rate dies too'':
$\varrho_\sigma$ is a guarantee, not a prediction. The exact transverse gap is the
subject of the next proposition, and \emph{it} can die for a different reason.)
\end{remark}

\begin{proposition}[gap collapse at a two-component limit]\label{prop:gap_collapse}
Suppose instead $p_t\to p_1$ where the limit chain has $m\geq2$ recurrent classes, and
the invariant measures converge, $\lambda(p_t)\to\lambda_1$ with $\lambda_1>0$
everywhere. Then the exact transverse Hessian gap vanishes:
$\zeta^{\perp}_{\min}\bigl(H(p_t)\bigr)\to0$ for
$H = g''(1)A^\dagger M_wA$, any fixed $w>0$.
\end{proposition}

\begin{proof}
With $m\geq2$ recurrent classes each charged by $\lambda_1$, the per-class invariant
densities $h_j = d\lambda^{(j)}/d\lambda_1$ span an $m$-dimensional kernel of $A_1 = P_1-I$
(density action), and some mean-zero combination $h\neq0$ has
$\langle h, H_1h\rangle=0$ with $\|h\|_{L^2(\lambda_1)}>0$. Since $\lambda_1$ is
everywhere positive, the quadratic forms $H(p_t)$ converge to $H_1$ on a
non-degenerating space, and the smallest transverse eigenvalue converges to $0$.
\end{proof}

\begin{remark}[the positivity hypothesis is necessary, measured]\label{rem:one_leak}
The hypothesis $\lambda_1>0$ everywhere requires the two leak rates --- into and out of
the would-be second class --- to close at comparable speeds, so that the limit measure is
a genuine mixture over both classes. A one-leak degeneration shows the hypothesis is
necessary: closing only the escape route drains the invariant measure onto the trap, the
limit operator \emph{on its own support} has no extra kernel, and the measured gap
stalls ($0.28\to0.014$ and no further on the one-leak path) --- while the two-leak path
takes it to $2.1\cdot10^{-9}$ (Figure~\ref{fig:app_path}, right; both paths are computed
by the released figure script). The distinction is exactly the
difference between ``the chain traps'' (hitting times diverge, $\varrho_\sigma$ dies
through $\sigma_*$) and ``the chain disconnects'' (the exact gap itself dies).
\end{remark}

\begin{remark}[minimum flow is adjacent to collapse]\label{rem:lp_degeneracy}
At a minimum-flow solution the LP leaves the backward rows of \emph{zero-flow states}
completely unconstrained: the min-flow face of the polytope contains kernels whose
unconstrained rows form traps. Flow regularization therefore steers toward a boundary
region where that trapping geometry is not a far-away pathology but a neighboring face.
The released architecture sweep reached it through a collapsed planner; a
flow-regularized trained $\mathcal P_B$ can reach it through the prices above ---
$\lamreg$ is the dial that chooses how close.
\end{remark}

\subsection{The entropic repair}\label{sec:entropic}

\begin{lemma}[trajectory-likelihood regularization]\label{lem:traj_entropy}
For interior-positive $\mathcal P_B$,
\begin{equation*}
    \sum_{s'} N(s')\,H\bigl(\mathcal P_B(\cdot\,|\,s')\bigr)
    \;=\; \mathbb E\bigl[-\log \mathcal P_B(\tau)\bigr],
\end{equation*}
the expected negative log-likelihood of the backward excursion --- so $N$-weighted
conditional entropy is \emph{path} entropy, exactly. (The $\lambda$-weighted variant
differs by the factor $2+\bar\sigma$; and neither is the edge-flow entropy
$\sum_e f_e\log f_e$ of entropic optimal transport, which differs further by the state
marginal term --- the relation to Sinkhorn is precise but not an identity, see
Remark~\ref{rem:entropic_ot}.)
\end{lemma}

\begin{proof}
Both sides are $\sum_{s'}\sum_u (\text{expected visits to } s')\,
\mathcal P_B(u|s')\allowbreak\,(-\log\mathcal P_B(u|s'))$, by Wald / the occupation
decomposition of the excursion.
\end{proof}

\begin{proposition}[soft-SSP keeps the constants alive]\label{prop:soft_ssp}
Regularize the slow objective with the relative entropy of the policy to the uniform
backward proposal, at temperature $\mu$:
\begin{equation*}
    \min_p\ \bar\sigma(p)\;+\;\mu\,\mathbb E_p\Bigl[\textstyle\sum_t
    \log\bigl(d_{x_t}\,p(u_t|x_t)\bigr)\Bigr]
\end{equation*}
(a pure path-entropy bonus --- the object of Lemma~\ref{lem:traj_entropy} --- differs
from this penalty by the trajectory's accumulated log-degrees; the KL-to-uniform form
is the one whose fixed point is stated below). Its optimizer is the soft-Bellman policy
$\pi_\mu(u|x)\propto e^{-V_\mu(u)/\mu}$ over parents, where $V_\mu$
solves $V_\mu(x) = 1-\mu\log\bigl(\tfrac1{d_x}\sum_{u\in\mathrm{pa}(x)}
e^{-V_\mu(u)/\mu}\bigr)$, $V_\mu(s_0)=0$. Then:
\begin{enumerate}
    \item \emph{(fixed point)} on any finite graph where $s_0$ is reachable from every
    state, value iteration from $0$ converges monotonically to a fixed point $V_\mu$
    with $V_\mu\leq(1+\mu\log d_{\max})\,V_{\mathrm{hard}}$;
    \item \emph{(policy floor)} $\pi_\mu(u|x)\geq e^{-\Delta_V/\mu}/d_{\max}$ with
    $\Delta_V$ the largest parent-value spread;
    \item \emph{(the certified floor)} wherever the minimal-occupation state is one soft
    choice off the hard path,
    \begin{equation*}
        \varrho_\sigma(\mu)\;\geq\;
        \frac{g''(1)\,e^{-\Delta_V/\mu}}{d_{\max}\,(2+\bar\sigma(\mu))\,\sigma_*(\mu)^2}
        \;>\;0 ;
    \end{equation*}
    in general $N(x)\geq r_{\min}\,(e^{-\Delta_V/\mu}/d_{\max})^{D}$ for $D$ the
    backward diameter and $r_{\min}$ the smallest positive sink-row mass, giving the
    same shape with a cruder exponent;
    \item \emph{(the price of the repair)} $\bar\sigma(\mu)-\bar\sigma_{\min} \leq
    \mu\,\log d_{\max}\,\bar\sigma(\mu)$.
\end{enumerate}
So the dial is explicit: the excess trajectory length is \emph{linear} in $\mu$ while
$\log\varrho_\sigma(\mu)$ is $\Theta(-1/\mu)$ --- the repair converts the deterministic
limit's exact zero (a void certificate) into a positive floor, at a cost in $\bar\sigma$
one can read off before training.
\end{proposition}

\begin{proof}
(1) The soft-Bellman operator $\mathsf T$ is monotone, and $\mathsf T0=\mathbf 1\geq0$,
so the iterates $V_n$ increase. The bound $V\leq c\,V_{\mathrm{hard}}$ with
$c:=1+\mu\log d_{\max}$ is invariant: using
$-\mu\log\bigl(\tfrac1d\sum e^{-V(u)/\mu}\bigr)\leq\min_uV(u)+\mu\log d$,
\begin{equation*}
    (\mathsf TV)(x)\;\leq\;1+c\min_{\mathrm{pa}(x)}V_{\mathrm{hard}}+\mu\log d_{\max}
    \;\leq\;c\,\bigl(1+\min_{\mathrm{pa}(x)}V_{\mathrm{hard}}\bigr)\;=\;c\,V_{\mathrm{hard}}(x),
\end{equation*}
since $c-1-\mu\log d_{\max}=0$ and $(c-1)\min V_{\mathrm{hard}}\geq0$; $V_{\mathrm{hard}}$
is finite by reachability. Monotone and bounded, $V_n$ converges pointwise; continuity of
$\mathsf T$ makes the limit a fixed point. (2) is immediate from the softmax form.
(4) $V_\mu$ is the optimal value of steps plus a \emph{non-negative} penalty (a
relative entropy per step), so $\mathbb E_{\pi_\mu}[\mathrm{steps}]\leq V_\mu$; and
$V_\mu\leq(1+\mu\log d_{\max})V_{\mathrm{hard}}$ by (1). Taking the sink-row
expectation, $\bar\sigma(\mu)\leq(1+\mu\log d_{\max})\bar\sigma_{\min}$, and
$\bar\sigma(\mu)-\bar\sigma_{\min}\leq\mu\log d_{\max}\,\bar\sigma_{\min}\leq
\mu\log d_{\max}\,\bar\sigma(\mu)$. (3) chains (2) through
$N_{\min}\geq\pi_\mu$-floor at the pivotal choice and
$\lambda_{\min}=N_{\min}/(2+\bar\sigma)$ into Proposition~\ref{prop:morozov_rate}(3);
the general form iterates the floor along a backward path of length at most $D$ from a
sink-row state of mass $\geq r_{\min}$.
\end{proof}

\begin{remark}[relation to entropic OT]\label{rem:entropic_ot}
Path-entropy regularization of a minimum-cost flow is the Schr\"odinger-bridge /
maximum-entropy-RL side of the entropic-OT family, not Sinkhorn itself: Sinkhorn
penalizes the coupling's entropy (edge flows), soft-SSP penalizes the \emph{policy}'s
path measure, and on graphs the two differ by the state-marginal entropy. What matters
here is not the taxonomy but the direction of the exchange: \emph{the certificate
constants of Section~\ref{sec:cv_divergence} are a quantitative reason to prefer the
entropic version of the flow regularizer} --- they are finite on the soft solution and
zero on the hard one, with the exchange rate of Proposition~\ref{prop:soft_ssp}.
\end{remark}

\emph{Summary of the bridge.} In the paper's variables, \eqref{equ:hse_objective}
reads: the loss holds the flow on the balanced manifold at $O(\lamreg^2)$
(Proposition~\ref{prop:two_timescale}); the regularizer walks the manifold by gradient
descent on $\bar\sigma$ (Lemma~\ref{lem:eq10_occupation}); the walk heads toward the
minimum-flow/OT boundary (Proposition~\ref{prop:drift_target}), where the certified
rate, the step cap and (at a disconnecting limit) the exact gap die at computable
speeds (Propositions~\ref{prop:constants_path} and~\ref{prop:gap_collapse}), next to
the min-flow face whose unconstrained backward rows form traps
(Remark~\ref{rem:lp_degeneracy}); and a path-entropy term restores every constant at a
price linear in the temperature (Proposition~\ref{prop:soft_ssp}).

\subsection{Numerical evidence}\label{sec:num_pb}

The script \texttt{audit\_pb\_spectral.py}, released with the code, computes everything
below exactly and exits nonzero on any violation.

\paragraph{The dictionary holds to machine precision.}
Over $12$ random loop-closed marked graphs: the occupation identity to
$4.3\cdot10^{-16}$; $\lambda(s_0)=\lambda(s_f)=1/(2+\bar\sigma)$ to $2.7\cdot10^{-16}$;
the interior sum $\sum\mathcal F/\mathcal F(s_f)=\bar\sigma$ and the full sum
$2+\bar\sigma$ of Lemma~\ref{lem:eq10_occupation} to $1.3\cdot10^{-13}$.

\paragraph{Minimum flow equals optimal transport, with a dual certificate.}
On a two-route instance and a four-terminal diamond: the minimum-interior-flow LP value
equals an independent transport-LP value (cost $=$ interior states on a shortest path)
to $10^{-10}$; complementary slackness against the solver's own dual potentials holds
with zero reduced cost on the support --- the tight-subgraph statement of
\citet[Thm~3.3]{maksimov2026ot} with the \emph{correct} potentials. (A fixed initial row
can force flow down a long route; tightness is relative to the optimal Kantorovich
potentials, not to naive minimum distances --- a naive min-distance version of this
check fails on exactly that instance.)

\paragraph{The spectral price is exact on the two-route instance.}
Along $\mathcal P_B(t)=(1-t)P_{\mathrm{unif}}+tP_{\mathrm{det}}$
(Figure~\ref{fig:app_path}, left): $N_{\min}=(1-t)/2$ \emph{exactly} ($10^{-12}$ at
every $t$); $\widehat B_\sigma$ grows as $(1-t)^{-1/2}$, from $12.65$ to $848.57$ at
$t=0.9999$; $\varrho_\sigma$ dies linearly, $1.25\cdot10^{-2}\to2.78\cdot10^{-6}$;
$\sigma_*$ \emph{decreases}, $4.000\to3.000$; $\|w\|_{L^\infty}$ grows to $1.1\cdot10^4$ for
uniform $\nutrain$ while $w\equiv1$ holds to $10^{-14}$ under adaptive replay
(Remark~\ref{rem:replay_rescue}). On the trap path with both leaks closing together
(Figure~\ref{fig:app_path}, right): the exact transverse gap falls
$2.84\cdot10^{-1}\to2.11\cdot10^{-9}$ and the limit kernel has exactly $2$ weak
invariant components, validating Proposition~\ref{prop:gap_collapse}; the certificate
ordering $\varrho_\sigma\leq\zeta^\perp_{\min}(H)$ holds at every point of both paths.

\paragraph{The soft-SSP floor is exactly tight.}
Fixed-point residual $<10^{-14}$ at every $\mu\in\{0.5,0.2,0.1,0.05,0.02\}$; the policy
floor and the certified $\varrho_\sigma$ floor of Proposition~\ref{prop:soft_ssp}(3)
hold at every $\mu$ --- and the floor is \emph{exactly tight} on this instance
($\varrho_\sigma$ equals it to three digits at every $\mu$, e.g.\ $9.670\cdot10^{-4}$ at
$\mu=0.5$ and $2.07\cdot10^{-45}$ at $\mu=0.02$); $\bar\sigma(\mu)$ decreases to the
hard value $2$ with excess inside the $\mu\log d_{\max}\,\bar\sigma$ bound at every
$\mu$ (Figure~\ref{fig:app_softssp}).

\begin{figure*}[tb]
\centering
\includegraphics[width=0.9\linewidth]{figures/figA5_path.pdf}
\caption{\textbf{Left:} the constants along the interpolation toward the minimum-flow
support (log scales): $\varrho_\sigma$ dies linearly with $N_{\min}=(1-t)/2$,
$\widehat B_\sigma$ diverges as $(1-t)^{-1/2}$, $\sigma_*$ stays bounded
(Proposition~\ref{prop:constants_path}). \textbf{Right:} the exact transverse gap on the
two-leak trap path against the one-leak stall (Proposition~\ref{prop:gap_collapse} and
Remark~\ref{rem:one_leak}).}
\label{fig:app_path}
\end{figure*}

\begin{figure*}[tb]
\centering
\includegraphics[width=0.9\linewidth]{figures/figA6_softssp.pdf}
\caption{The entropic repair's dial (Proposition~\ref{prop:soft_ssp}).
\textbf{Left:} $\varrho_\sigma(\mu)$ against its certified floor
$g''(1)e^{-\Delta_V/\mu}/(d_{\max}(2+\bar\sigma)\sigma_*^2)$ --- exactly tight on this
instance at every temperature. \textbf{Right:} $\bar\sigma(\mu)$'s excess over the hard
value against the linear bound $\mu\log d_{\max}\,\bar\sigma(\mu)$.}
\label{fig:app_softssp}
\end{figure*}
